\subsection{Example 2: the case of a free module}

Let \(k\) be a ring, \(\mathbf{M}\) a \(k\) -module, \(\mathbf{I}\) a set and \(p\) be an integer \(\geqslant 0\) . A map \(\boldsymbol{m}:1^{p}\to\mathbf{M}\) is said to be alternating if it satisfies the following two conditions:

\begin{enumerate}
\def\labelenumi{\alph{enumi})}
\tightlist
\item
  for every permutation \(\sigma \in \mathfrak{S}_p\) and every sequence \((\alpha_1, \ldots, \alpha_p) \in \mathrm{I}^p\) , we have
\end{enumerate}

\[
\boldsymbol {m} \left(\alpha_ {\sigma (1)}, \dots , \alpha_ {\sigma (p)}\right) = \varepsilon_ {\sigma} \boldsymbol {m} \left(\alpha_ {1}, \dots , \alpha_ {p}\right),
\]

\begin{enumerate}
\def\labelenumi{\alph{enumi})}
\setcounter{enumi}{1}
\tightlist
\item
  for every sequence \((\alpha_{1},\ldots,\alpha_{p})\in\mathbb{I}^{p}\) such that two of the indices \(\alpha_{1},\ldots,\alpha_{p}\) are equal, one has \(m(\alpha_{1},\ldots,\alpha_{p})=0\) .
\end{enumerate}

(In the case where I is a k-module and m is multilinear, one recovers the notion introduced in III, p.~80.)

Suppose I is finite and denote by \(C_{I}^{p}(M)\) the k-module of alternating maps from \(I^{p}\) to M.

Let \(L_{0}\) be a k-module, \((e_{i})_{i\in I}\) a family of elements of \(L_{0}\) ; we define two k-linear maps

\[
\begin{array}{l} g: \operatorname{Hom} _ {k} (\symbf {\Lambda} ^ {p} \mathrm{L} _ {0}, \mathrm{M}) \to \mathrm{C} _ {\mathrm{I}} ^ {p} (\mathrm{M}) \\ h: \mathrm{C} _ {\mathrm{I}} ^ {p} (\mathrm{M}) \to \mathrm{M} \otimes_ {k} \symbf {\Lambda} ^ {p} \mathrm{L} _ {0} \end{array}
\]

as follows: if \(f \in \operatorname{Hom}_k(\symbf{\Lambda}^p \mathbf{L}_0, \mathbf{M})\) , we set

\[
g (f) (\alpha_ {1}, \dots , \alpha_ {p}) = f (e _ {\alpha_ {1}} \wedge \dots \wedge e _ {\alpha_ {p}});
\]

let \(m \in C_{I}^{p}(M)\) , we define \(h(m) \in M \otimes_{k} \Lambda^{p} L_{0}\) . For every element \((\alpha_{1}, \ldots, \alpha_{p})\)

of \(I^{p}\) the element \(m(\alpha_{1},\ldots,\alpha_{p})\otimes(e_{\alpha_{1}}\wedge\ldots\wedge e_{\alpha_{p}})\) of \(\Lambda^{p}L_{0}\otimes_{k}M\) is zero if Card \(\{\alpha_{1},\ldots,\alpha_{p}\}<p\) and is independent of the order of the indices \(\alpha_{1},\ldots,\alpha_{p}\) if

\[
\text { Card } \{\alpha_ {1}, \dots , \alpha_ {p} \} = p.
\]

It depends only on the subset \(\mathbf{J} = (\alpha_{1},\ldots ,\alpha_{p})\) of I; denote it by \(h_\mathrm{J}(\pmb {m})\) ; we have \(h_\mathrm{J}(\pmb {m}) = 0\) if Card (J) \(< p\) ; we then set:

\[
h (\boldsymbol {m}) = \sum_ {\mathrm{J}} h _ {\mathrm{J}} (\boldsymbol {m}),
\]

where J ranges over the subsets of I having p elements.

Lemma 2. --- If the k-module \(\mathbf{L}_{0}\) is free with basis \((e_{i})_{i\in I}\) , the \(k\) -linear maps \(g\) and \(h\) are bijective.

This follows from III, p.~79, th. 1.

Now let M be a k-module, and \(\boldsymbol{x} = (x_{i})_{i \in I}\) a family of k-endomorphisms of M, pairwise commuting. Consider the polynomial ring \(\mathbf{A} = k[(\mathbf{X}_{i})_{i \in I}]\) and endow M with the A-module structure such that \(\mathrm{P}(\mathbf{X}_{i}) m = \mathrm{P}(x_{i}) m\) for \(P \in A\) and \(m \in M\) . Let, on the other hand, L be the free A-module \(A^{I}, (e_{i})_{i \in I}\) its canonical basis and \(u : L \to A\) the linear form which maps \(e_{i}\) over \(X_{i}\) for all \(i \in I\) . Consider the complexes of k-modules \(\mathbf{K}^{\mathbf{A}}(u, \mathbf{M})\) and \(\mathbf{K}_{\mathbf{A}}^{\bullet}(u, \mathbf{M})\) ; we have canonical isomorphisms

\[
\begin{array}{c} \mathbf {K} _ {\mathbf {A}} ^ {p} (u, \mathrm{M}) = \operatorname{Hom} _ {\mathbf {A}} \left(\boldsymbol {\Lambda} _ {\mathbf {A}} ^ {p} (\mathrm{A} ^ {\mathrm{I}}), \mathrm{M}\right) \to \operatorname{Hom} _ {k} \left(\boldsymbol {\Lambda} _ {k} ^ {p} (k ^ {\mathrm{I}}), \mathrm{M}\right), \\ \mathrm{M} \otimes_ {k} \boldsymbol {\Lambda} _ {k} ^ {p} (k ^ {\mathrm{I}}) \to \mathrm{M} \otimes_ {\mathbf {A}} \boldsymbol {\Lambda} _ {\mathbf {A}} ^ {p} (\mathrm{A} ^ {\mathrm{I}}) = \mathbf {K} _ {p} ^ {\mathbf {A}} (u, \mathrm{M}); \end{array}
\]

whence, by composition with the isomorphisms \(g\) and \(h\) , isomorphisms of \(k\) -modules

\[
\begin{array}{l} \theta^ {p}: \mathbf {K} _ {\mathrm{A}} ^ {p} (u, \mathrm{M}) \to \mathrm{C} _ {\mathrm{I}} ^ {p} (\mathrm{M}), \\ \theta_ {p}: \mathrm{C} _ {\mathrm{I}} ^ {p} (\mathrm{M}) \to \mathbf {K} _ {p} ^ {\mathrm{A}} (u, \mathrm{M}). \end{array}
\]

We denote by

\[
\begin{array}{l} \hat {c} ^ {p}: \mathbf {C} _ {\mathrm{I}} ^ {p} (\mathbf {M}) \to \mathbf {C} _ {\mathrm{I}} ^ {p + 1} (\mathbf {M}), \\ \hat {c} _ {p}: \mathbf {C} _ {\mathrm{I}} ^ {p} (\mathbf {M}) \to \mathbf {C} _ {\mathrm{I}} ^ {p - 1} (\mathbf {M}), \end{array}
\]

the \(k\) -homomorphisms obtained by transporting the differentials of \(\mathbf{K}_{\mathbf{A}}^{\bullet}(u, \mathbf{M})\) and \(\mathbf{K}_{\mathbf{A}}^{\mathbf{A}}(u, \mathbf{M})\) via the isomorphisms \(\theta\) . We have, for example:

\[
\left(\partial^ {p} \boldsymbol {m}\right) \left(\alpha_ {1}, \dots , \alpha_ {p + 1}\right) = \sum_ {j = 1} ^ {p + 1} (- 1) ^ {j + 1} x _ {\alpha_ {j}} \boldsymbol {m} \left(\alpha_ {1}, \dots , \alpha_ {j - 1}, \alpha_ {j + 1}, \dots , \alpha_ {p + 1}\right).\tag{12}
\]

The complex formed by the \(C_{I}^{p}(M)\) and the \(\partial^{p}\) (resp. the \(\partial_{p}\) ) is denoted \(\mathbf{K}^{\prime}(\mathbf{x},\mathbf{M})\) (resp. \(\mathbf{K}(\mathbf{x},\mathbf{M})\) ) and is called the ascending (resp. descending) Koszul complex

associated with the module M and the sequence of endomorphisms \((x_{1}, \ldots, x_{n})\) . We therefore have isomorphisms of complexes of k-modules

\[
\begin{array}{l} \theta^ {\bullet}: \mathbf {K} _ {\mathbf {A}} ^ {\bullet} (u, \mathbf {M}) \to \mathbf {K} ^ {\bullet} (x, \mathbf {M}), \\ \theta_ {\bullet}: \mathbf {K}. (x, \mathbf {M}) \to \mathbf {K} _ {\bullet} ^ {\mathbf {A}} (u, \mathbf {M}). \end{array}
\]

Remark. --- Conversely, let B be a \(k\) -algebra, L a free B-module with basis \((e_i)_{i \in I}\) , and M a B-module. The data of a linear form \(u: L \to B\) is equivalent to that of a family \(x = (x_i)_{i \in 1}\) of elements of B, via the relation \(x_i = u(e_i)\) . The complex of \(k\) -modules underlying \(\mathbf{K}_{\mathrm{B}}^{\bullet}(u, \mathrm{M})\) (resp. \(\mathbf{K}^{\mathrm{B}}(u, \mathrm{M})\) ) is then identified, via the isomorphism \(\theta^{\bullet}\) (resp. \(\theta\) .) with the Koszul complex \(\mathbf{K}'(x, \mathrm{M})\) (resp. \(\mathbf{K}^{\bullet}(x, \mathrm{M})\) ). For example \(\mathbf{K}^{\mathrm{B}}(u)\) is identified with \(\mathbf{K}'(x, \mathrm{B})\) .

We introduce, as in no. 1 (X, p.~147), the notations H.(x, M), H.(x, M), etc., and all the results of nos. 1 and 2 apply mutatis mutandis, the module A \(^{I}\) being free. For example, we have isomorphisms

\[
\begin{array}{l} \mathrm{H} _ {0} (\boldsymbol {x}, \mathrm{M}) \to \mathrm{M} / (\boldsymbol {x}) \mathrm{M} \\ \mathrm{H} ^ {0} (\boldsymbol {x}, \mathrm{M}) \to \mathrm{Hom} _ {\mathrm{A}} \left(\mathrm{A} / (\boldsymbol {x}), \mathrm{M}\right), \end{array}
\]

where \((\pmb{x})\) denotes the ideal \(\sum \mathbf{A}x_{i}\) of A. Also for example, if \(\mathbf{K}.(\pmb{x},\mathbf{A})\) is acyclic in degrees \(>0\) , we have isomorphisms

\[
\begin{array}{l}\mathrm{H} _ {r} (\boldsymbol {x}, \mathrm{M}) \rightarrow \operatorname{Tor} _ {r} ^ {\mathrm{A}} (k, \mathrm{M}),\\\operatorname{Ext} _ {\mathrm{A}} ^ {r} (k, \mathrm{M}) \rightarrow \mathrm{H} ^ {r} (\boldsymbol {x}, \mathrm{M}).\end{array}
\]

Finally, suppose I is (finite and) totally ordered, for example \(I = \{1, \ldots, n\}\) ; let us identify \(\symbf{\Lambda}^{n}(\mathbf{A}^{\mathrm{I}})\) to A by means of the basis element \(e_{1} \wedge \ldots \wedge e_{n}\) and let us translate the isomorphism \(\mathbf{K}_{p}^{\mathrm{A}}(u, \mathrm{M}) \to \mathbf{K}_{\mathrm{A}}^{n-p}(u, \mathrm{M})\) of X, p.~149. It becomes, by transport by \((\theta_{p})\) and \((\theta^{n-p})\) , the isomorphism

\[
\mathrm{C} _ {\mathrm{I}} ^ {p} (\mathbf {M}) \rightarrow \mathrm{C} _ {\mathrm{I}} ^ {n - p} (\mathbf {M})
\]

which associates with \(m \in C_{I}^{p}(M)\) the element \(\check{m}\) of \(C_{1}^{n-p}(M)\) such that

\[
\boldsymbol {m} \left(\alpha_ {1}, \dots , \alpha_ {p}\right) = \check {\boldsymbol {m}} \left(\beta_ {1}, \dots , \beta_ {n - p}\right)\tag{13}
\]

if \((\alpha_{1},\ldots,\alpha_{p},\beta_{1},\ldots,\beta_{n-p})\) is an even permutation of \(\{1,\ldots,n\}\) Let us also note that when \(I=\{1,\ldots,n\}\) , one can identify \(C_{I}^{p}(M)\) with the set of families \(m(\alpha_{1},\ldots,\alpha_{p})\) of elements of M where \(\alpha_{1}<\alpha_{2}<\ldots<\alpha_{p}\) ; formula (12) remains valid, as does relation (13).

Example. --- Take \(\mathbf{M} = k[\mathbf{T}_1, \ldots, \mathbf{T}_n]\) ; the Koszul complex \(\mathbf{K}\) ( \(\partial/\partial\mathbf{T}\) , \(\mathbf{M}\) ) associated with the sequence of endomorphisms ( \(\partial/\partial\mathbf{T}_1, \ldots, \partial/\partial\mathbf{T}_n\) ) is identified with the de Rham complex of \(k[x_{1},\ldots,x_{n}]\) over \(k\) (X, p.~44): to \(\boldsymbol{m}\in\mathbf{C}_{\{1,\ldots,n\}}^{p}(\mathbf{M})\) , one associates the differential form

\[
\omega (\boldsymbol {m}) = \sum_ {\alpha_ {1} <   \dots <   \alpha_ {p}} \boldsymbol {m} \left(\alpha_ {1}, \dots , \alpha_ {p}\right) d x _ {\alpha_ {1}} \wedge \dots \wedge d x _ {\alpha_ {p}},
\]

cf.~formula (12) and example 1, p.~44.

